\section{Método}
\label{desenvolvimento}

Trata-se de uma pesquisa experimental: os fatores (modelo, precisão numérica, resolução e taxa da câmera) são manipulados e as métricas de desempenho e de qualidade são medidas sob as mesmas condições. Todo o código, as versões fixadas das ferramentas e os dados brutos estão no repositório do experimento\footnote{\url{https://github.com/joaoleaogf/deteccao-pessoas-borda}}, e cada etapa é refeita com um único comando.

\subsection{Ambiente experimental}

Sem acesso ao hardware físico (Raspberry Pi 4/5 e acelerador Coral Edge TPU), o dispositivo de borda é representado por um contêiner Docker \texttt{linux/arm64} executado de forma nativa em uma máquina virtual ARM64 da Oracle Cloud (VM.Standard.A1.Flex, processador Ampere com núcleos Neoverse-N1), com 2 núcleos e 12\,GB de memória. Uma primeira tentativa por emulação (QEMU em x86-64) foi abandonada: sob o emulador, a biblioteca Python do LiteRT mostrou-se cerca de dez vezes mais lenta que a ferramenta nativa de \textit{benchmark}, o que tornaria a coleta inviável no prazo. As medições parciais emuladas ficaram no repositório apenas como anexo.

\subsection{Modelos e quantização}

Foram avaliados três detectores pré-treinados no COCO: YOLOv8n \citep{jocher2023yolo}, SSD-MobileNet~V1 \citep{liu2016ssd,howard2017mobilenets} e EfficientDet-Lite0 \citep{tan2020efficientdet}, cada um em FP32 e INT8, no formato do LiteRT \citep{litert2026}. O INT8 do SSD-MobileNet~V1 e do EfficientDet-Lite0 é inteiro, inclusive entrada e saída, obtido por quantização pós-treinamento \citep{jacob2018quantization} calibrada nas mesmas 128 imagens (COCO128). O INT8 do YOLOv8n é o exportado oficialmente pela Ultralytics, que mantém entrada, saída e cabeça de detecção em ponto flutuante (INT8 híbrido). Os modelos INT8 também foram submetidos ao Edge TPU Compiler, versão 14.1 \citep{coral2026compiler}.

\subsection{Fatores e amostragem}

Os fatores são o modelo (3 níveis), a precisão (FP32 e INT8), a resolução do quadro da câmera (320, 480 e 640 pixels) e a taxa da câmera (5, 15 e 30 FPS), o que dá 54 cenários na borda. A resolução é a do quadro entregue pela câmera: o YOLOv8n usa o modelo exportado na mesma resolução, enquanto o SSD-MobileNet~V1 (entrada fixa de $300\times300$) e o EfficientDet-Lite0 ($320\times320$) redimensionam o quadro, e esse custo entra no pré-processamento.

Por restrição de tempo, o desenho original (1.050 inferências em cada cenário) foi reduzido: cada uma das 18 combinações modelo $\times$ precisão $\times$ resolução foi medida uma única vez, em processo novo, com 350 quadros, dos quais os 50 primeiros foram descartados como aquecimento (\nQuadrosValidos{} quadros válidos). Como a taxa da câmera não altera o tempo de processamento de cada quadro, os cenários de 5, 15 e 30 FPS foram derivados dessas medições: o quadro $k$ fica disponível em $k/\text{fps}$, o processamento começa quando o quadro está disponível e o anterior terminou, e a espera é registrada como tempo de fila, separado da latência fim a fim. Também se conta a fração de quadros que uma política de descarte (processar sempre o quadro mais recente) deixaria de processar.

\subsection{Conjunto de teste e métricas de qualidade}

O conjunto de teste tem 350 imagens do COCO val2017 \citep{lin2014coco} com pessoas, sem anotações de multidão, balanceadas entre 1, 2 e 3 ou mais pessoas e ordenadas pelo identificador; a lista e o resumo criptográfico (sha256) de cada imagem estão no repositório. Cada imagem é redimensionada para a resolução da câmera e entregue como quadro. A qualidade é medida pelo mAP@50 da classe pessoa \citep{cocoapi2026} e pelo F1 com limiar de pontuação 0,5 e IoU 0,5, a partir das detecções salvas na própria coleta. A decodificação das saídas e a supressão de não máximos (IoU 0,5, até 100 detecções por quadro) são feitas com os mesmos parâmetros para os três modelos e foram conferidas em x86 antes da coleta.

\subsection{Métricas de desempenho}

Para cada quadro são registrados o tempo de pré-processamento, o de inferência (2 \textit{threads}) e o de pós-processamento, cuja soma é a latência fim a fim, além do uso de CPU do processo e da memória residente. A taxa sustentável é o inverso da latência média. A energia por quadro seria estimada pelo produto entre a potência de referência da placa e a latência, mas não foi calculada por falta de um valor de potência com fonte verificável.

\subsection{Baseline de nuvem}

Para a QP2, o mesmo conjunto de quadros é enviado a um servidor de inferência: dois contêineres \texttt{linux/amd64} em uma rede Docker, com o enlace limitado por \texttt{tc netem} a um tempo de ida e volta (RTT) de 40\,ms e subida de 10\,Mbit/s. Esses dois valores são assumidos. \pendente{justificar RTT e banda de subida com uma fonte} O cliente apenas codifica cada quadro em JPEG (qualidade 80) e o envia por HTTP; o servidor executa os mesmos modelos FP32. São registrados os tempos de codificação, transmissão, inferência e retorno e o tamanho de cada quadro, de onde se calcula a banda de subida necessária em cada taxa da câmera. Para manter o pareamento por quadro, usou-se envio quadro a quadro em JPEG no lugar de RTSP ou WebRTC.

\subsection{Análise estatística}

Para cada cenário calculam-se média, desvio padrão, mediana, percentil 95 e intervalo de confiança de 95\% da média. As comparações FP32 $\times$ INT8 são pareadas por quadro: aplica-se o teste de Shapiro-Wilk às diferenças e, conforme o resultado, o teste $t$ pareado ou o de Wilcoxon, com $\alpha = 0{,}05$.
